We study a pricing game in markets with budget-constrained buyers, where each seller offers a single good at a price she sets, and each buyer chooses a utility-maximizing bundle subject to her budget. We consider both the rivalrous setting, such as markets for physical goods, and the non-rivalrous setting, such as markets for datasets. Pure Nash equilibria often fail to exist in this game. However, Nash equilibrium assumes that the other sellers keep their prices fixed when one seller deviates, whereas in practice, sellers harmed by a price change may respond, for example, by matching the new price.

We propose a new equilibrium notion called \emph{Response-Aware Rational Equilibrium} (RARE). At a RARE, every deviation by a seller can be deterred by a response from the other sellers, in which no responder earns less than she did right after the deviation. Every Nash equilibrium is a RARE.
Assuming that buyers' purchases (among their optimal choices) change continuously with prices, we show that a RARE always exists for rivalrous goods, and for non-rivalrous goods when buyers have equal budgets. Moreover, the RARE that we construct extracts the buyers' entire budget. For non-rivalrous goods with unequal budgets, we show that an approximate RARE exists, where a deviating seller can increase her revenue by a factor of at most $b_{\mathrm{avg}}/b_{\min}$, the ratio of the average budget to the minimum budget. Our proof uses two nested applications of the KKM lemma.
We also study a natural dynamic in which sellers repeatedly raise their prices to earn the revenue that they could guarantee themselves by deviating, and show that it reaches an approximate RARE, with additive error $\delta$, within $\lfloor B/\delta \rfloor$ rounds, where $B$ is the buyers' total budget.
The main results of this paper are formally verified in the Lean 4 proof assistant
\texorpdfstring{\cite{eklavya2026lean}}{(https://github.com/sharmaeklavya2/dm-rare)}.
